% GENERATED FILE. Edit canonical lessons, then run npm run generate:books.
% canonical-source-hash: fnv1a64:06a85e263a021676
% canonical-lessons: TA-C486-tatavu, TA-C486-amukku, TA-C486-koppali, TA-C486-unrukol, TA-C486-mil

\chapter{Care at Home}
\label{ch:ta-care-at-home}
\hlchaptermodality{486}

\begin{chapteropening}
\textbf{By the end of this chapter:} \emph{I can talk about rubbing on ointment, massaging, gargling and recovering.}

Complete the last lesson of chapter 486: I can talk about rubbing on ointment, massaging, gargling and recovering.
\end{chapteropening}

\section[\texorpdfstring{\ta{தடவு} (taṭavu)}{தடவு (taṭavu)}]{\ta{தடவு} (taṭavu) — to rub on; to apply}

\label{lesson:TA-C486-tatavu}



\begin{quote}
Before the new one: say the Tamil for a medicated oil, then the Tamil for an ointment.
\end{quote}

\subsection*{You'll want to know: \ta{தடவு}}
\textbf{\ta{தடவு}} — \emph{taṭavu} — \textquotedblleft{}to rub on; to apply\textquotedblright{}.

\begin{quote}
\ta{காயத்தில்} \ta{அம்மா} \ta{களிம்பு} \ta{தடவினாள்}.
\end{quote}

\begin{quote}
\emph{kāyattil ammā kaḷimpu taṭaviṉāḷ.}
\end{quote}

\begin{quote}
Mother put ointment on the wound.
\end{quote}

Past (I …): \ta{தடவினேன்}, \emph{taṭaviṉēṉ}.

Also to stroke, or to grope one’s way in the dark.

\begin{cousinweb}
\textbf{\ta{தடவு}} (\emph{taṭavu}) is native Dravidian (DEDR 3025), of stroking and of smearing on oil.
\end{cousinweb}

\begin{grammarlens}[title={What you've built}]
One more word for describing what you lived and felt.
\end{grammarlens}

\subsection*{Your turn}
\begin{itemize}
\raggedright
  \item \emph{Say it:} \emph{taṭavu}
  \item \emph{Say it:} \emph{taṭavu}, once more
  \item \emph{Say it:} \emph{kāyattil ammā kaḷimpu taṭaviṉāḷ.} — Mother put ointment on the wound.
  \item \emph{Recall:} say the Tamil for a medicated oil, then the Tamil for an ointment, then say \emph{taṭavu} again
\end{itemize}

\begin{tcolorbox}[breakable,colback=teal!4,colframe=teal!35!black,title={Before you move on}]
What does \ta{தடவு} mean? (\textquotedblleft{}To rub on; to apply\textquotedblright{}.) Say it once more.
\end{tcolorbox}
\section[\texorpdfstring{\ta{அமுக்கு} (amukku)}{அமுக்கு (amukku)}]{\ta{அமுக்கு} (amukku) — to press down; to massage}

\label{lesson:TA-C486-amukku}



\begin{quote}
Before the new one: say the Tamil for an ointment, then the Tamil for to rub on.
\end{quote}

\subsection*{You'll want to know: \ta{அமுக்கு}}
\textbf{\ta{அமுக்கு}} — \emph{amukku} — \textquotedblleft{}to press down; to massage\textquotedblright{}.

\begin{quote}
\ta{தாத்தாவின்} \ta{கால்களை} \ta{நான்} \ta{அமுக்கினேன்}.
\end{quote}

\begin{quote}
\emph{tāttāviṉ kālkaḷai nāṉ amukkiṉēṉ.}
\end{quote}

\begin{quote}
I massaged grandfather’s legs.
\end{quote}

Past (I …): \ta{அமுக்கினேன்}, \emph{amukkiṉēṉ}.

The strong partner of \ta{அமுங்கு} \textquotedblleft{}to sink, be pressed down\textquotedblright{}.

\begin{cousinweb}
\textbf{\ta{அமுக்கு}} (\emph{amukku}) is native Dravidian (DEDR 169), the transitive partner of \ta{அமுங்கு} \emph{amuṅku}, \textquotedblleft{}to be pressed down\textquotedblright{}.
\end{cousinweb}

\begin{grammarlens}[title={What you've built}]
One more word for describing what you lived and felt.
\end{grammarlens}

\subsection*{Your turn}
\begin{itemize}
\raggedright
  \item \emph{Say it:} \emph{amukku}
  \item \emph{Say it:} \emph{amukku}, once more
  \item \emph{Say it:} \emph{tāttāviṉ kālkaḷai nāṉ amukkiṉēṉ.} — I massaged grandfather’s legs.
  \item \emph{Recall:} say the Tamil for an ointment, then the Tamil for to rub on, then say \emph{amukku} again
\end{itemize}

\begin{tcolorbox}[breakable,colback=teal!4,colframe=teal!35!black,title={Before you move on}]
What does \ta{அமுக்கு} mean? (\textquotedblleft{}To press down; to massage\textquotedblright{}.) Say it once more.
\end{tcolorbox}
\section[\texorpdfstring{\ta{கொப்பளி} (koppaḷi)}{கொப்பளி (koppaḷi)}]{\ta{கொப்பளி} (koppaḷi) — to gargle}

\label{lesson:TA-C486-koppali}



\begin{quote}
Before the new one: say the Tamil for to rub on, then the Tamil for to press down.
\end{quote}

\subsection*{You'll want to know: \ta{கொப்பளி}}
\textbf{\ta{கொப்பளி}} — \emph{koppaḷi} — \textquotedblleft{}to gargle\textquotedblright{}.

\begin{quote}
\ta{தொண்டை} \ta{வலிக்கு} \ta{உப்புத்} \ta{தண்ணீரில்} \ta{வாய்} \ta{கொப்பளித்தேன்}.
\end{quote}

\begin{quote}
\emph{toṇṭai valikku upput taṇṇīril vāy koppaḷittēṉ.}
\end{quote}

\begin{quote}
For my sore throat I gargled with salt water.
\end{quote}

Past (I …): \ta{கொப்பளித்தேன்}, \emph{koppaḷittēṉ}.

\ta{வாய்} \ta{கொப்பளி}, to rinse the mouth.

\begin{cousinweb}
\textbf{\ta{கொப்பளி}} (\emph{koppaḷi}) is native Dravidian (DEDR 2107). DEDR keeps it apart from \ta{கொப்புளம்} \textquotedblleft{}blister\textquotedblright{} (2106).
\end{cousinweb}

\begin{grammarlens}[title={What you've built}]
One more word for describing what you lived and felt.
\end{grammarlens}

\subsection*{Your turn}
\begin{itemize}
\raggedright
  \item \emph{Say it:} \emph{koppaḷi}
  \item \emph{Say it:} \emph{koppaḷi}, once more
  \item \emph{Say it:} \emph{toṇṭai valikku upput taṇṇīril vāy koppaḷittēṉ.} — For my sore throat I gargled with salt water.
  \item \emph{Recall:} say the Tamil for to rub on, then the Tamil for to press down, then say \emph{koppaḷi} again
\end{itemize}

\begin{tcolorbox}[breakable,colback=teal!4,colframe=teal!35!black,title={Before you move on}]
What does \ta{கொப்பளி} mean? (\textquotedblleft{}To gargle\textquotedblright{}.) Say it once more.
\end{tcolorbox}
\section[\texorpdfstring{\ta{ஊன்றுகோல்} (ūṉṟukōl)}{ஊன்றுகோல் (ūṉṟukōl)}]{\ta{ஊன்றுகோல்} (ūṉṟukōl) — a walking stick; a crutch}

\label{lesson:TA-C486-unrukol}



\begin{quote}
Before the new one: say the Tamil for to press down, then the Tamil for to gargle.
\end{quote}

\subsection*{You'll want to know: \ta{ஊன்றுகோல்}}
\textbf{\ta{ஊன்றுகோல்}} — \emph{ūṉṟukōl} — \textquotedblleft{}a walking stick; a crutch\textquotedblright{}.

\begin{quote}
\ta{கால்} \ta{உடைந்ததால்} \ta{அவர்} \ta{ஊன்றுகோல்} \ta{வைத்து} \ta{நடந்தார்}.
\end{quote}

\begin{quote}
\emph{kāl uṭaintatāl avar ūṉṟukōl vaittu naṭantār.}
\end{quote}

\begin{quote}
Because his leg was broken, he walked with a crutch.
\end{quote}

\ta{ஊன்று} \textquotedblleft{}to plant firmly\textquotedblright{} + \emph{kōl} \textquotedblleft{}a stick\textquotedblright{}: a stick you lean on.

\begin{cousinweb}
\textbf{\ta{ஊன்றுகோல்}} (\emph{ūṉṟukōl}) is a native compound of \ta{ஊன்று} \emph{ūṉṟu}, \textquotedblleft{}to plant firmly, to lean upon\textquotedblright{} (DEDR 763), and \emph{kōl}, \textquotedblleft{}a stick\textquotedblright{}.
\end{cousinweb}

\begin{grammarlens}[title={What you've built}]
One more word for describing what you lived and felt.
\end{grammarlens}

\subsection*{Your turn}
\begin{itemize}
\raggedright
  \item \emph{Say it:} \emph{ūṉṟukōl}
  \item \emph{Say it:} \emph{ūṉṟukōl}, once more
  \item \emph{Say it:} \emph{kāl uṭaintatāl avar ūṉṟukōl vaittu naṭantār.} — Because his leg was broken, he walked with a crutch.
  \item \emph{Recall:} say the Tamil for to press down, then the Tamil for to gargle, then say \emph{ūṉṟukōl} again
\end{itemize}

\begin{tcolorbox}[breakable,colback=teal!4,colframe=teal!35!black,title={Before you move on}]
What does \ta{ஊன்றுகோல்} mean? (\textquotedblleft{}A walking stick; a crutch\textquotedblright{}.) Say it once more.
\end{tcolorbox}
\section[\texorpdfstring{\ta{மீள்} (mīḷ)}{மீள் (mīḷ)}]{\ta{மீள்} (mīḷ) — to pull through; to recover}

\label{lesson:TA-C486-mil}



\begin{quote}
Before the new one: say the Tamil for to gargle, then the Tamil for a walking stick.
\end{quote}

\subsection*{You'll want to know: \ta{மீள்}}
\textbf{\ta{மீள்}} — \emph{mīḷ} — \textquotedblleft{}to pull through; to recover\textquotedblright{}.

\begin{quote}
\ta{அவர்} \ta{நோயிலிருந்து} \ta{மெதுவாக} \ta{மீண்டார்}.
\end{quote}

\begin{quote}
\emph{avar nōyiliruntu metuvāka mīṇṭār.}
\end{quote}

\begin{quote}
He slowly recovered from the illness.
\end{quote}

Past (I …): \ta{மீண்டேன்}, \emph{mīṇṭēṉ}.

From an illness or a shock: -\ta{இலிருந்து} \ta{மீள்}. \ta{மீண்டும்} \textquotedblleft{}again\textquotedblright{} shares the root.

\begin{cousinweb}
\textbf{\ta{மீள்}} (\emph{mīḷ}) is native Dravidian (DEDR 4883): to return, to be cured, to be rescued.
\end{cousinweb}

\begin{grammarlens}[title={What you've built}]
One more word for describing what you lived and felt.
\end{grammarlens}

\subsection*{Your turn}
\begin{itemize}
\raggedright
  \item \emph{Say it:} \emph{mīḷ}
  \item \emph{Say it:} \emph{mīḷ}, once more
  \item \emph{Say it:} \emph{avar nōyiliruntu metuvāka mīṇṭār.} — He slowly recovered from the illness.
  \item \emph{Recall:} say the Tamil for to gargle, then the Tamil for a walking stick, then say \emph{mīḷ} again
\end{itemize}

\begin{tcolorbox}[breakable,colback=teal!4,colframe=teal!35!black,title={Before you move on}]
What does \ta{மீள்} mean? (\textquotedblleft{}To pull through; to recover\textquotedblright{}.) Say it once more.
\end{tcolorbox}
